\section{The investor's problem}
\label{sec:05_investor}

The investor chooses among a sale, zero, and a purchase.
Each class compares one correctly signed order with zero once we exclude the wrong sign.
The key bounds apply under every candidate profile, before investor optimality is imposed. \paper{Sec.~3.3; Lemma~A.3}

\subsection{Hold the price and entry rules fixed}
\label{sec:05_investor-fixed}

A candidate satisfies (E1)--(E4) and (E6); it need not satisfy investor optimality (E5).
An investor deviation changes only its own order.
It faces the candidate's price function $P(\cdot)$ and preparation rule $\zeta(\cdot)$, as the caution in \Cref{sec:02_model-equilibrium} explains.
Since its order fills the whole slot, a deviation to $q$ reaches $P(q)$ and $e(q)$.
Do not recompute Bayes beliefs from the deviating order profile: those beliefs belong to a different candidate.

Conditional on investor arrival, the expected payoff in class $j$ is
\[
U_j(q)=q\{t_0+e(q)(t_j-t_0)-P(q)\}-k|q|,
\qquad q\in\{-1,0,1\}.
\]
The factor $\lambda>0$ of the ex ante payoff changes no comparison between orders.

\subsection{Derive the residual advantage and the six payoffs}
\label{sec:05_investor-residual}

Use $w_j=t_j-t_0$, $w_H-w_L=\DT$, the conditional value $\E[V_T\mid j,X=q]=t_0+e(q)w_j$ and the price
$P(q)=t_0+e(q)[w_L+\DT\mu(q)]$ of \Cref{sec:04_slot_beliefs-price}.
In class $H$, subtract the price from the conditional value:
\begin{align*}
\E[V_T\mid H,X=q]-P(q)
&=t_0+e(q)w_H-t_0-e(q)[w_L+\DT\mu(q)]\\
&=e(q)[w_H-w_L-\DT\mu(q)]\\
&=e(q)\DT[1-\mu(q)].
\end{align*}
In class $L$, a sale earns the opposite difference:
\begin{align*}
P(q)-\E[V_T\mid L,X=q]
&=t_0+e(q)[w_L+\DT\mu(q)]-t_0-e(q)w_L\\
&=e(q)\DT\mu(q).
\end{align*}
The proceeds without entry, $t_0$, cancel in both expressions.
The common entry gain $e(q)w_L$ also cancels.

Multiply the residual by the order sign and subtract the cost.
The six payoffs are
\begin{align*}
U_H(1)&=e(1)\DT[1-\mu(1)]-k,& U_H(0)&=0,\\
U_H(-1)&=-e(-1)\DT[1-\mu(-1)]-k,\\[2pt]
U_L(-1)&=e(-1)\DT\mu(-1)-k,& U_L(0)&=0,\\
U_L(1)&=-e(1)\DT\mu(1)-k.
\end{align*}
Define the gross returns of the two correctly signed nonzero orders:
\[
K_H=e(1)\DT[1-\mu(1)],\qquad
K_L=e(-1)\DT\mu(-1).
\]
Then a correct order has net return $K_j-k$.

\begin{intuition}[What rational prices remove and what they leave]
Market makers anticipate preparation and include its expected effect in the price.
The investor therefore earns no extra return from that common expected effect alone.
It still knows the class while market makers retain some uncertainty.
Its residual advantage is entry times the target-payoff spread times that residual uncertainty.
A liquidity order prevents either nonzero order from revealing the class with certainty.
\end{intuition}

This argument also applies to general class distributions.
The investor knows the class, rather than the realized acquisition value.
The relevant spread is $t_H-t_L$, the difference between expected target proceeds in the two classes.

\subsection{Wrong-signed orders are dominated}
\label{sec:05_investor-sign}

Under the class ranking, $\DT\ge0$.
Also $e(x)\ge0$ and $0\le\mu(x)\le1$.
Thus
\[
U_H(-1)\le-k<0=U_H(0),\qquad
U_L(1)\le-k<0=U_L(0).
\]
Zero strictly beats a sale in class $H$ and a purchase in class $L$.
No equilibrium puts positive probability on either wrong-signed order.
This conclusion comes from payoffs; it is not a restriction on the candidate profiles.

Class $H$ therefore compares $K_H$ with $k$.
It buys if $K_H>k$ and abstains if $K_H<k$.
Class $L$ sells if $K_L>k$ and abstains if $K_L<k$.
At equality, the correct order and zero are both best responses.
Rule IM permits any mixture of these best responses; rule I0 selects zero.

\subsection{The gross-return bounds under every candidate}
\label{sec:05_investor-bounds}

For a correct purchase, the posterior bound gives $\mu(1)\le M(\pi)$.
Subtract from one, multiply by $\DT\ge0$, and then use the entry floor:
\begin{align*}
1-\mu(1)&\ge1-M(\pi),\\
\DT[1-\mu(1)]&\ge\DT[1-M(\pi)],\\
K_H=e(1)\DT[1-\mu(1)]&\ge\rho\DT[1-\mu(1)]
\ge\rho\DT[1-M(\pi)].
\end{align*}
For a correct sale, $\mu(-1)\ge m(\pi)$ gives
\[
K_L=e(-1)\DT\mu(-1)\ge\rho\DT\mu(-1)\ge\rho\DT m(\pi).
\]
For the upper bounds, $e\le1$, $\mu\le1$, and $1-\mu\le1$ give $K_H,K_L\le\DT$.
Set $\kappa(\pi)=\min\{m(\pi),1-M(\pi)\}$.
We obtain
\begin{resultbox}[The bound that fixes orders before a candidate is known]
Under every candidate profile,
\[
\rho\kappa(\pi)\DT\le K_H\le\DT,\qquad
\rho\kappa(\pi)\DT\le K_L\le\DT.
\]
These are bounds on the gross returns of correctly signed nonzero orders. \paper{(A.6)}
They hold for every $\beta<1/2$.
\end{resultbox}

\begin{caution}[Scope of the lower bound]
Do not apply the positive lower bound to zero or to a wrong-signed order.
Zero earns zero, and a wrong-signed order has a negative net return.
Do not replace $\kappa$ by $m$ at an arbitrary prior.
The two coincide at $\pi=1/2$.
\end{caution}

A mixed candidate that is not an equilibrium still obeys these bounds.
The bounds use only Bayes' rule and the entry floor, not the optimality of the orders.
\Cref{sec:benchmark-investor} gives a worked candidate.

\subsection{Why the prior one half gives the largest floor}
\label{sec:05_investor-prior}

Calculate the residual at the top posterior:
\begin{align*}
1-M(\pi)
&=\frac{1-\pi}{1+\pi u}\\
&=\frac{1-\pi}{(1-\pi)+\pi(1+u)}=m(1-\pi).
\end{align*}
The function $m$ rises strictly with its argument.
Indeed,
\[
m(\pi)=\frac{\pi}{1+u-\pi u},\qquad
m'(\pi)=\frac{1+u}{(1+u-\pi u)^2}>0.
\]
It follows that
\[
\kappa(\pi)=m(\min\{\pi,1-\pi\}),\qquad
\kappa(1-\pi)=\kappa(\pi).
\]
The minimum of $\pi$ and $1-\pi$ is largest at $1/2$.
Therefore $\kappa$ has its unique maximum there:
\[
\kappa(1/2)=\frac1{2+u}.
\]
When $\pi>1/2$, the bound for class $H$ binds because the market already favors $H$.
When $\pi<1/2$, the bound for class $L$ binds for the same reason.
Away from the prior one half, the floor falls on both sides by the same amount.

\subsection{Why the weak economy forces zero and the strong economy forces full orders}
\label{sec:05_investor-force}

If $\DT<k$, every correct nonzero order satisfies
\[
U_j(q)=K_j-k\le\DT-k<0=U_j(0).
\]
Wrong-signed orders are also strictly worse than zero.
Thus zero is the unique best response of both classes under every candidate, and (E5) excludes every mixture.

If $k<\rho\kappa\DT$, every correct order instead satisfies
\[
U_j(q)=K_j-k\ge\rho\kappa\DT-k>0=U_j(0).
\]
The wrong sign is strictly worse than zero.
Thus each class has one unique best response: a purchase in $H$, a sale in $L$.
Every equilibrium uses full orders; the argument excludes all other pure and mixed profiles at once.

The weak economy of Proposition~2 has a spread below the cost, so it falls in the first case.
The strong economy has a floor above the cost, so it falls in the second case.
\Cref{sec:benchmark-investor} gives the benchmark margins.

\begin{caution}[Uniqueness and verification are two tests]
The bound $\rho\kappa\DT$ holds against every candidate, so it proves uniqueness.
The returns $K_H$ and $K_L$ at the full-order tuple verify that one candidate only.
Keep the two tests apart on the panel, as in the old guide.
\end{caution}

\begin{panelnote}[What changed from the Laplace version]
The slot has only unit orders and zero, so the proof compares three payoffs and needs no derivative in the order size.
The Laplace factor $1-1/s$ (the old guide's $1-1/b$) therefore drops out of the trading-cost window.
\end{panelnote}
